\chapter{Những sự thật về thất bại}
\markboth{Những sự thật về thất bại}{Truths About Failure}


\section{Thất bại đến bất ngờ}
\begin{truyen}{Thất bại đến bất ngờ}{Life lesson}
\chuhoa{A}{nh nghĩ mọi thứ ổn. Một ngày dự án đổ vỡ, đối tác rút. Thất bại đến bất ngờ.}
\textit{(He thought everything was fine. One day the project collapsed, the partner left. Failure comes unexpectedly.)}
\end{truyen}

\begin{baihoc}
Đừng tưởng thành công là vĩnh viễn.
\textit{(Don't assume success is forever.)}
\end{baihoc}

\begin{ghinhoanh}
\textbf{Short English to remember:}

Don't assume success is forever.

\end{ghinhoanh}

\begin{tuhocnhanh}
1. Câu chuyện này gợi cho bạn suy nghĩ gì? \textit{What does this story make you think?}

2. Bạn sẽ nhớ bài học nào nhất? \textit{Which lesson will you remember most?}
\end{tuhocnhanh}
\ngancach


\section{Ai cũng từng sai}
\begin{truyen}{Ai cũng từng sai}{Life lesson}
\chuhoa{K}{hông ai anh ngưỡng mộ mà chưa từng vấp. Ai cũng từng sai — chỉ khác cách họ đứng dậy.}
\textit{(No one he admired had never stumbled. Everyone has been wrong — the difference is how they get up.)}
\end{truyen}

\begin{baihoc}
Sai không xấu; không chịu sửa mới xấu.
\textit{(Being wrong isn't bad; refusing to fix it is.)}
\end{baihoc}

\begin{ghinhoanh}
\textbf{Short English to remember:}

Being wrong isn't bad; refusing to fix it is.

\end{ghinhoanh}

\begin{tuhocnhanh}
1. Câu chuyện này gợi cho bạn suy nghĩ gì? \textit{What does this story make you think?}

2. Bạn sẽ nhớ bài học nào nhất? \textit{Which lesson will you remember most?}
\end{tuhocnhanh}
\ngancach


\section{Thất bại làm con người tỉnh táo}
\begin{truyen}{Thất bại làm con người tỉnh táo}{Life lesson}
\chuhoa{S}{au lần phá sản, anh mới nhìn rõ ai thật ai giả, mình mạnh chỗ nào yếu chỗ nào. Thất bại làm con người tỉnh táo.}
\textit{(After bankruptcy he saw clearly who was real and who wasn't, where he was strong and weak. Failure makes you clear-headed.)}
\end{truyen}

\begin{baihoc}
Thất bại là bài học đắt nhưng sâu.
\textit{(Failure is an expensive but deep lesson.)}
\end{baihoc}

\begin{ghinhoanh}
\textbf{Short English to remember:}

Failure is an expensive but deep lesson.

\end{ghinhoanh}

\begin{tuhocnhanh}
1. Câu chuyện này gợi cho bạn suy nghĩ gì? \textit{What does this story make you think?}

2. Bạn sẽ nhớ bài học nào nhất? \textit{Which lesson will you remember most?}
\end{tuhocnhanh}
\ngancach


\section{Người mạnh là người đứng dậy}
\begin{truyen}{Người mạnh là người đứng dậy}{Life lesson}
\chuhoa{N}{hiều người ngã giống nhau. Khác nhau ở chỗ ai đứng dậy, ai nằm mãi. Người mạnh là người đứng dậy.}
\textit{(Many fall the same way. The difference is who gets up and who stays down. The strong one is the one who gets up.)}
\end{truyen}

\begin{baihoc}
Sức mạnh không nằm ở không ngã.
\textit{(Strength isn't in never falling.)}
\end{baihoc}

\begin{ghinhoanh}
\textbf{Short English to remember:}

Strength isn't in never falling.

\end{ghinhoanh}

\begin{tuhocnhanh}
1. Câu chuyện này gợi cho bạn suy nghĩ gì? \textit{What does this story make you think?}

2. Bạn sẽ nhớ bài học nào nhất? \textit{Which lesson will you remember most?}
\end{tuhocnhanh}
\ngancach


\section{Sai lầm là thầy giáo tốt}
\begin{truyen}{Sai lầm là thầy giáo tốt}{Life lesson}
\chuhoa{A}{nh học từ sai lầm nhiều hơn từ thành công. Mỗi lần sai, anh biết thêm một đường không nên đi. Sai lầm là thầy giáo tốt.}
\textit{(He learned more from mistakes than from success. Each error showed him a path not to take. Mistakes are good teachers.)}
\end{truyen}

\begin{baihoc}
Đừng sợ sai — sợ không học từ sai.
\textit{(Don't fear being wrong — fear not learning from it.)}
\end{baihoc}

\begin{ghinhoanh}
\textbf{Short English to remember:}

Don't fear being wrong — fear not learning from it.

\end{ghinhoanh}

\begin{tuhocnhanh}
1. Câu chuyện này gợi cho bạn suy nghĩ gì? \textit{What does this story make you think?}

2. Bạn sẽ nhớ bài học nào nhất? \textit{Which lesson will you remember most?}
\end{tuhocnhanh}
\ngancach


\section{Người thành công từng thất bại nhiều lần}
\begin{truyen}{Người thành công từng thất bại nhiều lần}{Life lesson}
\chuhoa{A}{nh đọc tiểu sử những người nổi tiếng — họ đều có chuỗi thất bại. Người thành công từng thất bại nhiều lần.}
\textit{(He read the bios of famous people — they all had a string of failures. Successful people have failed many times.)}
\end{truyen}

\begin{baihoc}
Thất bại là bước đệm, không phải điểm dừng.
\textit{(Failure is a stepping stone, not a stop.)}
\end{baihoc}

\begin{ghinhoanh}
\textbf{Short English to remember:}

Failure is a stepping stone, not a stop.

\end{ghinhoanh}

\begin{tuhocnhanh}
1. Câu chuyện này gợi cho bạn suy nghĩ gì? \textit{What does this story make you think?}

2. Bạn sẽ nhớ bài học nào nhất? \textit{Which lesson will you remember most?}
\end{tuhocnhanh}
\ngancach


\section{Không ai đi thẳng đến thành công}
\begin{truyen}{Không ai đi thẳng đến thành công}{Life lesson}
\chuhoa{C}{on đường anh đi quanh co, có đoạn tối. Không ai đi thẳng đến thành công.}
\textit{(The path he took was winding, with dark stretches. No one goes straight to success.)}
\end{truyen}

\begin{baihoc}
Chấp nhận đường dài — đó mới là thật.
\textit{(Accept the long road — that's what's real.)}
\end{baihoc}

\begin{ghinhoanh}
\textbf{Short English to remember:}

Accept the long road — that's what's real.

\end{ghinhoanh}

\begin{tuhocnhanh}
1. Câu chuyện này gợi cho bạn suy nghĩ gì? \textit{What does this story make you think?}

2. Bạn sẽ nhớ bài học nào nhất? \textit{Which lesson will you remember most?}
\end{tuhocnhanh}
\ngancach


\section{Thất bại giúp mình hiểu bản thân}
\begin{truyen}{Thất bại giúp mình hiểu bản thân}{Life lesson}
\chuhoa{K}{hi mọi thứ sụp đổ, anh mới thấy mình có thể chịu đựng đến đâu, cần gì thật sự. Thất bại giúp mình hiểu bản thân.}
\textit{(When everything collapsed he saw how much he could endure and what he really needed. Failure helps you understand yourself.)}
\end{truyen}

\begin{baihoc}
Trong đau, mình thấy rõ mình hơn.
\textit{(In pain, you see yourself more clearly.)}
\end{baihoc}

\begin{ghinhoanh}
\textbf{Short English to remember:}

In pain, you see yourself more clearly.

\end{ghinhoanh}

\begin{tuhocnhanh}
1. Câu chuyện này gợi cho bạn suy nghĩ gì? \textit{What does this story make you think?}

2. Bạn sẽ nhớ bài học nào nhất? \textit{Which lesson will you remember most?}
\end{tuhocnhanh}
\ngancach


\section{Nỗi đau dạy những bài học lớn}
\begin{truyen}{Nỗi đau dạy những bài học lớn}{Life lesson}
\chuhoa{A}{nh không quên những đêm mất ngủ, những lần nước mắt. Nỗi đau dạy những bài học lớn.}
\textit{(He doesn't forget the sleepless nights, the tears. Pain teaches the biggest lessons.)}
\end{truyen}

\begin{baihoc}
Đau không vô ích nếu mình chịu học.
\textit{(Pain isn't useless if you're willing to learn.)}
\end{baihoc}

\begin{ghinhoanh}
\textbf{Short English to remember:}

Pain isn't useless if you're willing to learn.

\end{ghinhoanh}

\begin{tuhocnhanh}
1. Câu chuyện này gợi cho bạn suy nghĩ gì? \textit{What does this story make you think?}

2. Bạn sẽ nhớ bài học nào nhất? \textit{Which lesson will you remember most?}
\end{tuhocnhanh}
\ngancach


\section{Sau thất bại vẫn còn con đường}
\begin{truyen}{Sau thất bại vẫn còn con đường}{Life lesson}
\chuhoa{A}{nh từng nghĩ mình hết cửa. Năm năm sau anh đứng ở nơi khác. Sau thất bại vẫn còn con đường.}
\textit{(He once thought he had no way out. Five years later he stood somewhere else. After failure there is still a path.)}
\end{truyen}

\begin{baihoc}
Đừng dừng ở thất bại — bước tiếp.
\textit{(Don't stop at failure — keep walking.)}
\end{baihoc}

\begin{ghinhoanh}
\textbf{Short English to remember:}

Don't stop at failure — keep walking.

\end{ghinhoanh}

\begin{tuhocnhanh}
1. Câu chuyện này gợi cho bạn suy nghĩ gì? \textit{What does this story make you think?}

2. Bạn sẽ nhớ bài học nào nhất? \textit{Which lesson will you remember most?}
\end{tuhocnhanh}
\ngancach
